\documentclass[10pt]{article}
\usepackage[T1]{fontenc}
\usepackage[utf8]{inputenc}
\usepackage[french]{babel}
\usepackage{amsmath,amssymb,mathrsfs,mathpazo}
\usepackage{geometry}
\geometry{a4paper,margin=2cm}
\usepackage{graphicx,array,tabularx,tabularray,colortbl,multirow,multicol,booktabs,enumitem,pifont,ulem,fancybox,diagbox,cancel}
\usepackage{pstricks,pstricks-add,pst-plot,pst-tree,pst-node,pst-text,pst-func,pst-math,pst-bspline,pst-3dplot,pst-eucl,pst-solides3d}
\let\pscalculate\undefined
\usepackage{pst-bezier}
\usepackage{tikz,pgfplots,tkz-tab}
\usetikzlibrary{arrows,arrows.meta,patterns,calc,positioning,trees,decorations.markings,shapes.geometric,shapes.misc,shapes.arrows,intersections,decorations.pathmorphing,decorations.pathreplacing,backgrounds,shadows,babel}
\pgfplotsset{compat=1.18}
\usepgfplotslibrary{fillbetween}
\usepackage[np]{numprint}
\usepackage{listings,makecell,pas-tableur,twoopt}
\let\barre\undefined
\usepackage[pstricks]{tablvar}
\definecolor{darkgreen}{RGB}{0,100,0}
\providecommand{\red}{\color{red}}
\providecommand{\blue}{\color{blue}}
\providecommand{\green}{\color{green}}
\providecommand{\black}{\color{black}}
\providecommand{\white}{\color{white}}
\providecommand{\gray}{\color{gray}}
\providecommand{\cyan}{\color{cyan}}
\providecommand{\magenta}{\color{magenta}}
\providecommand{\yellow}{\color{yellow}}
\providecommand{\violet}{\color{violet}}
\providecommand{\orange}{\color{orange}}
\providecommand{\euro}{\texteuro}
\providecommand{\EUR}{\texteuro}
\providecommand{\ieme}{\textsuperscript{e}}
\providecommand{\degres}{\ensuremath{^\circ}}
\providecommand{\degree}{\ensuremath{^\circ}}
\providecommand{\vv}[1]{\overrightarrow{#1}}
\providecommand{\Coord}[1]{#1}
\providecommand{\decosix}{}
\providecommand{\starredbullet}{\ensuremath{\bullet}}
\pagestyle{empty}
\setlength{\parindent}{0pt}
\usepackage{adjustbox}\DeclareUnicodeCharacter{2236}{:}\DeclareUnicodeCharacter{2190}{\ensuremath{\leftarrow}}
\providecommand{\R}{\mathbb{R}}
\providecommand{\N}{\mathbb{N}}
\providecommand{\Z}{\mathbb{Z}}
\providecommand{\Q}{\mathbb{Q}}
\providecommand{\D}{\mathbb{D}}
\providecommand{\C}{\mathbb{C}}
\providecommand{\e}{{\mathrm{e}}}
\providecommand{\vect}[1]{\overrightarrow{#1}}
\providecommand{\vectt}[1]{\overrightarrow{\mathrm{#1}}}
\providecommand{\vvt}[1]{\overrightarrow{\mathrm{#1}}}
\providecommand{\np}[2][]{{#2~#1}}
\providecommand{\vv}[1]{{\overrightarrow{#1}}}
\renewcommand{\d}{{\mathrm{d}}}
\providecommand{\degres}{{^\circ}}
\providecommand{\degre}{{^\circ}}
\providecommand{\blue}{\color{blue}}
\providecommand{\red}{\color{red}}
\providecommand{\enskip}{\quad}
\providecommand{\llbracket}{\lbrack\!\lbrack}
\providecommand{\rrbracket}{\rbrack\!\rbrack}
\providecommand{\euro}{{\text{€}}}
\providecommand{\no}{{\text{n°}}}
\providecommand{\interval}[2]{{[#1;#2]}}
\providecommand{\intervFO}[2]{{[#1;#2[}}
\providecommand{\dotfill}{\ldots}

\providecommand{\ds}{\displaystyle}
\providecommand{\textding}[1]{#1}
\providecommand{\Oij}{\(\mathrm{O};\vec{\imath},\vec{\jmath}\)}
\providecommand{\Oijk}{\(\mathrm{O};\vec{\imath},\vec{\jmath},\vec{k}\)}
\providecommand{\Ouv}{\(\mathrm{O};\vec{u},\vec{v}\)}
\providecommand{\intervFF}[2]{\ensuremath{[#1;#2]}}
\renewcommand{\euro}{\ifmmode\text{\texteuro}\else\texteuro\fi}
\DeclareUnicodeCharacter{00B0}{\ifmmode^{\circ}\else\textdegree\fi}
\DeclareUnicodeCharacter{20AC}{\ifmmode\text{\texteuro}\else\texteuro\fi}
% Énoncés : Éducation nationale. Transcription et corrigés : APMEP, Denis Vergès et contributeurs. Découpage et classement : CoopMaths, Sébastien Lozano, Éric Elter et contributeurs. Adaptation de la présentation : Maths Explorer.
% CC BY-SA : https://coopmaths.fr/www/about/
% https://coopmaths.fr/alea/static/bac/2025/tex/bac_2025_09_sujet2_metropole_4.tex
\begin{document}
\noindent\textit{APMEP / CoopMaths. CC BY-SA. Présentation adaptée par Maths Explorer.}\par\medskip
\medskip

L'espace est rapporté à un repère orthonormé \Oijk.

On considère les points 

\begin{center}A$(4~;~-1~;~3)$,\quad B$(-1~;~1~;~-2)$,\quad C(0~;~4~;~5) et D$(-3~;~-4~;~6)$.\end{center}

\begin{enumerate}
\item
	\begin{enumerate}
		\item Vérifier que les points A, B, C ne sont pas alignés.
	\end{enumerate}

On admet qu’une équation cartésienne du plan (ABC) est : $29x + 30y - 17z = 35$.

	\begin{enumerate}[resume]
		\item Les points A, B, C, D sont-ils coplanaires ? Justifier.
	\end{enumerate}
\end{enumerate}

On admet que lorsque quatre points ne sont pas coplanaires, il existe un unique point situé à égale distance de ces quatre points.

L’objectif de cet exercice est de déterminer le point H se situant à égale distance des quatre points A, B, C, D.

On définit le plan médiateur d’un segment comme le plan passant par le milieu de ce segment et orthogonal à la droite portant ce segment. C’est l’ensemble des points équidistants des extrémités de ce segment.

\begin{enumerate}[resume]
\item Soit $P_1$ le plan médiateur du segment [AB].
	\begin{enumerate}
		\item Déterminer les coordonnées du milieu du segment [AB].
		\item En déduire qu'une équation cartésienne de $P_1$ est : $5x - 2y + 5z = 10$.
	\end{enumerate}
\item On note $P_2$ le plan médiateur du segment [CD].
	\begin{enumerate}
		\item Soit M un point du plan $P_2$ de coordonnées $(x~;~y~;~z)$.
		
Exprimer MC$^2$ et MD$^2$ en fonction des coordonnées de M.

En déduire qu’une équation cartésienne du plan $P$ est : $-3x - 8y + z = 10$.
		\item Justifier que les plans $P_1$ et $P_2$ sont sécants.
	\end{enumerate}
\item Soit $\Delta$ la droite dont une représentation paramétrique est :
\[\left\{\begin{array}{l c l}
x& =& -2 - 1,9t\\
y&= &\phantom{-2 - 1,9}t\\
z&= & \phantom{-}4 + 2,3t
\end{array}\right.\: \text{où}\: t \in \R\]

Démontrer que $\Delta$ est la droite d’intersection de $P_1$ et $P_2$.
\end{enumerate}

On note $P_3$ le plan médiateur du segment [AC].

On admet qu'une équation cartésienne du plan $P_3$ est : $8x - 10y - 4z= - 15.$

\begin{enumerate}[resume]
\item Démontrer que la droite $\Delta$ et le plan $P_3$ sont sécants.
\item Justifier que le point d’intersection entre $\Delta$ et $P_3$ est le point H.
\end{enumerate}

\end{document}
